
Giản đồ không-thời gian là một công cụ giúp ta hình dung trực quan chuyển động của các vật thể (coi là các chất điểm) trong không gian và theo thời gian, tương ứng với trục toạ độ (trục hoành) và trục thời gian (trục tung). Thông tin về một sự kiện bao gồm vị trí và thời điểm nó xảy ra được xác định bởi một điểm trên giản đồ, và việc một chất điểm nằm ở vị trí nào đó tại thời điểm nào đó cũng chính là một sự kiện. Do đó, chuyển động của chất điểm vạch ra trên giản đồ không-thời gian một đường, gọi là đường thế giới.
\vspace{8pt}

Một hệ quy chiếu được xác định bởi một tập hợp điểm không chuyển động so với nhau, trong hệ quy chiếu này, các điểm ấy vạch ra các đường (thế giới) thẳng đứng. Các đường t=const nằm ngang đan xen với các đường trên lập thành một lưới toạ độ mà trên đó kích thước các ô được chuẩn hoá. Như vậy không-thời gian của một hệ quy chiếu biểu diễn trên giản đồ chính là một hệ lưới, thay đổi hệ quy chiếu là thay đổi hệ lưới đó. 
\vspace{8pt}

Trong cơ học cổ điển, thời gian độc lập với mọi quan sát viên, thành thử các đường nằm ngang trong hệ lưới này vẫn nằm ngang trong hệ lưới khác. Đó là thời gian tuyệt đối. 

\subsection*{Tương đối tính Galilei}
Xét hệ ba vật: cây (đứng yên), xe (vận tốc $v=10 \mathrm{m/s}$) và bóng (vận tốc $u=5 \mathrm{m/s}$), chuyển động cùng chiều với xe. Tại $t=0$, tọa độ cây và xe trùng nhau, quả bóng cách cây một khoảng $10\,\mathrm{m}$. Các thông số trên được xác định trong hệ quy chiếu mặt đất. Mô hình hóa chuyển động dưới dạng một chiều, các vật được coi là chất điểm. (xem \figref{P1_tikz1}).
\vspace{8pt}

\begin{figure}[htbp]
\centering
\begin{tikzpicture}[x=0.48cm,y=0.58cm,>=Latex,line cap=round,line join=round]
  \tikzset{
    axis/.style={->,very thick,black},
    originmark/.style={circle,fill=red!85!black,inner sep=1.7pt},
    carcolor/.style={draw=magenta!85!black},
    ball/.style={circle,fill=blue,draw=blue,inner sep=2.1pt},
    helper/.style={thin,black!35}
  }

  \draw[axis] (0,0) -- (24,0) node[below right] {$x\,(\mathrm{m})$};
  \node[originmark] at (0,0) {};
  \node[above,text=black] at (0,0) {$O \equiv C$};
  \coordinate (C0) at (0,0);
  \node[ball,label={[above,text=blue]$B$}] (B0) at (10,0) {};
  \draw[<->,helper] (0,-0.8) -- (10,-0.8);
  \node[below,text=black] at (5,-0.8) {$10\,\mathrm{m}$};
  \draw[->,very thick,magenta!85!black] (C0.center) ++(0,0) -- ++(5,0)
    node[above,text=magenta!85!black] {$v_C=10\,\mathrm{m/s}$};
  \draw[->,very thick,blue] (B0.center) ++(0,0) -- ++(2.5,0)
    node[below,text=blue] {$v_B=5\,\mathrm{m/s}$};
\end{tikzpicture}
\caption{Hình minh họa}
\label{P1_tikz1}
\end{figure}

\begin{enumerate}[label=(\alph*)]
\item Thiết lập các đường thế giới của ba vật và vẽ giản đồ không--thời gian cho hệ vật trên trong hai hệ quy chiếu gắn với cây $\left( x, t \right)$ và gắn với xe $\left( x', t' \right)$. Hình vẽ thể hiện rõ số liệu đề cho. 

\begin{figure}[htbp]
\centering
\begin{tikzpicture}[x=1cm,y=1cm,>=Latex,line cap=round,line join=round,font=\sansmath\sffamily]
  \tikzset{
    axis/.style={->,line width=1pt,black},
    basegrid/.style={gray!55,line width=1pt,dash pattern=on 4pt off 3pt},
    movinggrid/.style={magenta!28,line width=1pt,dash pattern=on 6pt off 3pt},
    worldcar/.style={violet!90!black,line width=1pt},
    worldball/.style={blue!90!black,line width=1pt}
  }

  \def\xmin{-2}
  \def\xmax{6}
  \def\ymax{5}

  % background square grid
  \draw[basegrid] (\xmin,0) grid (\xmax,\ymax);

  % moving-frame oblique grid: all lines are parallel to the car worldline
  % and clipped to stay strictly inside the square grid frame.
  \begin{scope}
    \clip (\xmin,0) rectangle (\xmax,\ymax);
    % 2 lines on the left, 5 lines on the right of the main purple line
    \foreach \x in {-1,0,1,2,3,4,5} {
      \draw[movinggrid] (\x,0) -- (\x+5,5);
    }
    \foreach \y in {0,1,2,3,4} {
      \draw[movinggrid] (\xmin,\y) -- (\xmin+5,\y+5);
    }
  \end{scope}

  % axes
  \draw[axis] (\xmin,0) -- (7,0) node[below right=2pt] {$x\equiv x'$};
  \draw[axis] (0,0) -- (0,6) node[above left=1pt] {$t$};

  % worldlines
  \draw[worldcar] (0,0) -- (5,5);
  \draw[axis] (5,5) -- (6,6) node[above right=1pt] {$t'$};
  \draw[worldball] (1,0) -- (3,5);

  % origin mark
  \fill[red!85!black] (0,0) circle (2pt);
  \node[below left=1pt,text=red!85!black] at (0,0) {$O$};

  % distance annotation: one square = 10 m
  \draw[gray!75,<->,line width=1pt] (0,-1) -- (1,-1);
  \draw[gray!75,line width=1pt] (0,-1) -- (0,0);
  \draw[gray!75,line width=1pt] (1,-1) -- (1,0);
  \node[below=2pt,text=red!85!black] at (0.5,-1) {$10\,\mathrm{m}$};
\end{tikzpicture}
\caption{Minh họa sự chồng chập xét trong hệ quy chiếu cây}
\label{P1_tikz2}
\end{figure}

\item Vì hai hệ quy chiếu này chỉ là hai cách mô tả cùng một diễn biến, ta có thể chồng chập các hệ lưới tọa độ của chúng lên cùng một giản đồ thay vì vẽ riêng rẽ. Khi chọn hệ quy chiếu gắn với cây làm gốc, đường thế giới của xe không còn là đường thẳng đứng mà trở thành một đường xiên; do đó, các đường lưới của hệ gắn với xe cũng nghiêng đi so với hệ gốc. Biểu diễn này cho phép quan sát sự biến đổi của hệ lưới mà không cần sử dụng các phép biến đổi đại số phức tạp. \Figref{P1_tikz2} minh họa giản đồ chồng chập lấy hệ quy chiếu gắn cây làm gốc. Biểu diễn hệ lưới của cây, các đường thế giới trong giản đồ của xe. Dựa vào hình dạng của các ô lưới mới, xác định tọa độ $\left( x, t \right)$ của một điểm khi biết tọa độ $\left( x', t' \right)$ của nó. Nêu rõ lập luận trong các bước vẽ, vẽ đúng tỉ lệ.
\end{enumerate}

\subsection*{Tương đối tính Einstein}
Cơ học Galilei có giới hạn của nó: thực nghiệm cho thấy tốc độ ánh sáng trong chân không luôn bằng trong mọi hệ quy chiếu quán tính -- bất kể hệ đó đang chuyển động thế nào. Với vận tốc $v$ đáng kể so với $c$, việc khảo sát sẽ làm xuất hiện các sai lệch mà cơ học thông thường không giải thích được. Đây là lúc ta dùng mô hình tương đối tính của Einstein.
\vspace{8pt}

Giản đồ Minkowski là một mô hình hình học trực quan giúp biểu diễn chuyển động tương đối của các vật có vận tốc đáng kể so với ánh sáng. Để thuận tiện cho việc khảo sát, cũng như đồng nhất không--thời gian, ta sử dụng hệ tọa độ $(x, ct)$. Khi khảo sát sự kiện trong hai hệ quy chiếu khác nhau, vì thực tại quan sát được là duy nhất, các sự kiện vật lý phải có sự liên hệ chặt chẽ thông qua các phép biến đổi hình học và toán học.
\vspace{8pt}

Xét một cái cây đứng yên trên mặt đất, và một chiếc ô tô đang chuyển động sang bên phải với vận tốc bằng một nửa tốc độ ánh sáng (xem \figref{P1_I3}). Cả cái cây và chiếc xe đều có đồng hồ, các đồng hồ được chỉnh để trùng nhau tại thời điểm $t=t'=0$ khi cả hai cùng ở cùng một vị trí với một bóng đèn. Bóng đèn và đồng hồ được bật vào lúc xe bắt đầu chuyển động. Khi đèn bật, hai photon (hạt ánh sáng) sẽ được phát ra từ bóng đèn với tốc độ ánh sáng $c$ ngược chiều nhau. Ta xét bài toán chuyển động một chiều, mô hình hóa các vật thể như các chất điểm.

\begin{enumerate}[label=(\alph*)]
\addtocounter{enumi}{2}
    \item Chỉ ra sự mâu thuẫn với nguyên lý bất biến tốc độ ánh sáng của giả định $t=t'$; chứng minh rằng các sự kiện đồng thời trong hệ này, về tổng quát, sẽ không còn đồng thời trong hệ kia. Khi ấy, liệu tập hợp các sự kiện xảy ra đồng thời trong hệ quy chiếu này có còn lập thành một đường thẳng trong một hệ quy chiếu khác không? Nếu có, liệu trục $O\mathrm{x}$ có còn trùng với trục $O\mathrm{x}'$ không?
    \item Xác hệ lưới của xe trong hệ quy chiếu của cây và ngược lại. Gọi $\theta$ là góc hợp giữa 2 trục $O\mathrm{x}$ và $O\mathrm{x}'$. Hãy lập luận chỉ ra $\theta$ cũng là góc hợp giữa hai trục $c t$ và $c t'$. Tổng quát hóa bài toán, chỉ ra rằng $\tan \theta = v/c = \beta$, với v là vận tốc của xe đối với đất. Biện luận các giới hạn của $\tan \theta$.
    
\begin{figure}[htbp]
\centering
% 1D model on Ox axis only (ignore y, z directions)
\definecolor{lightaccent}{HTML}{FF5858}
\definecolor{caraccent}{HTML}{D33EFF}
\begin{tikzpicture}[font=\sansmath\sffamily, x=0.48cm,y=0.98cm,>=Latex,line cap=round,line join=round]
  \tikzset{
    axis/.style={->,very thick,black},
    tree/.style={circle,fill=blue!70!black,inner sep=2.1pt},
    car/.style={rectangle,fill=caraccent,minimum width=8pt,minimum height=6pt},
    helper/.style={thin,black!35}
  }

  \draw[axis] (0,0) -- (24,0) node[below right] {$x\,(\mathrm{m})$};
  \fill (0,0) circle (2pt);
  \node[below,text=black] at (0,0) {$O \equiv C$};
  \coordinate (C0) at (0,0);
  \draw[->,very thick,lightaccent] (C0.center) ++(0,0) -- ++(6,0)
    node[above] {$\vec c$};
  \draw[->,very thick,caraccent] (C0.center) ++(0,0) -- ++(3,0)
    node[above] {$\vec v_C=0.5c$};
  \draw[->,very thick,lightaccent] (C0.center) ++(0,0) -- ++(-6,0)
    node[above] {$\vec c$};
\end{tikzpicture}
\caption{Hình minh họa}
\label{P1_I3}
\end{figure}

\newpage

    \item Một tính chất quan trọng của không--thời gian là: diện tích của mỗi ô lưới đơn vị là bảo toàn khi chuyển hệ quy chiếu. Áp dụng tính chất đó, xác định tỷ lệ giữa một giây trên trục $c t$ và trục $c t'$; tỷ lệ giữa một mét trên trục $x$ và trục $x'$.
    \item Ba hiện tượng nổi bật trong thuyết tương đối hẹp bao gồm: Sự mất đồng thời của các sự kiện, Sự giãn thời gian, và Sự co độ dài. Dựa vào các tính chất của giản đồ không-thời gian mới được giới thiệu, ta sẽ chứng minh các hệ quả quan trọng này.
    \begin{enumerate}[label=(\roman*)]
        \item Sự giãn thời gian: Khi một vật chuyển động so với người quan sát, các khoảng thời gian đo bởi đồng hồ gắn với vật đó sẽ trông dài hơn khi nhìn từ hệ khác. Xét một khoảng thời gian đơn vị $\Delta t'$ trên đồng hồ gắn với xe. Nhìn từ cây, khoảng thời gian trên là $\Delta t$. Chứng minh hệ thức liên hệ: $$\Delta t = \gamma \Delta t' $$
        \item Sự co độ dài: Khi một vật chuyển động so với người quan sát, chiều dài của nó đo theo phương chuyển động sẽ trông ngắn lại khi nhìn từ hệ khác. Xét một cây thước có chiều dài riêng $L_0$ nằm yên trong hệ cây. Đo từ xe, cây thước có chiều dài $L$. Chứng minh hệ thức liên hệ: $$L = \gamma L_0$$
    \end{enumerate}
Cho biết $\gamma$ là một hệ số chỉ phụ thuộc vào $v$ và $c$.
\end{enumerate}
\vspace{8pt}